\par\Needspace{7\baselineskip}\bigskip
\begin{center}
\pdfbookmark[2]{Gauss to Encke, Göttingen, 25 March 1818}{pallas-brief-18180325}
\textsc{Gauss} to \textsc{Encke}, Göttingen, 25 March 1818.
\end{center}
\noindent\hspace*{\parindent}. . . . .\quad I venture once more to call upon your readiness, so often demonstrated, to oblige me with your great skill in computation. I hope that the subject will also be of some interest to you yourself. My table for the perturbations of Pallas is now quite complete. Its ease of use leaves me nothing to desire. With the aid of the table I have compared all 12 oppositions observed so far, from 1803 . . . . 1817, with the elements once again, refined the latter, and obtained splendid agreement. In doing so I made no use at all of the numerous observations from 1802, which only began after the opposition, and I am very eager to know how closely my elements agree with the observations from that year. It may yet be possible to derive the opposition from them, since the comparison can be made with elements that must come so close to the truth that \SourceLeaf{437}even observations somewhat farther removed can be used. Would you therefore be so good as
\begin{quote}
»to compare precisely all the meridian observations\\
»of Pallas from April and May 1802 with the\\
»osculating elements?«
\end{quote}
\par My tables give the following elements. The node and perihelion are referred to the mean equinox of the epoch and are to be regarded as fixed relative to the stars.
\begin{center}
\begingroup\small\setlength{\tabcolsep}{4pt}
\begin{tabular}{@{}lr@{}}
\multicolumn{2}{c}{1802 April 19,49214 Göttingen time}\\[3pt]
Mean longitude . . . . . . . . . . . . . .&\(166^0\;44'\;30''{,}43\)\\
Perihelion . . . . . . . . . . . . . . . . .&\(121\;12\;17{,}28\)\\
Ascending node . . . . . . . .&\(172\;27\;35{,}19\)\\
Angle of eccentricity . . . . . . . . . .&\(14\;13\;27{,}83\)\\
Inclination of the orbit . . . . . . . . . . .&\(34\;37\;47{,}88\)\\
Daily mean tropical motion . . . .&\(770''{,}47396\)\\
Logarithm of the semimajor axis . . . .&\(0{,}442\,2173.\)
\end{tabular}
\endgroup
\end{center}
\noindent You may first correct the epoch of longitude so that\quad . . . . .
\par
